% =====================================================================
%  Terms with more than one meaning (back matter)
% =====================================================================
\chapter*{Terms with more than one meaning}\label{ch:glossary}
\phantomsection
\addcontentsline{toc}{chapter}{Terms with more than one meaning}
\markboth{Terms with more than one meaning}{Terms with more than one meaning}

Several words are used in this book with more than one meaning, because different fields chose the
same word for different ideas. The list gives the meanings and the sections where each one is used.

\begin{description}
\item[Bias] (1) The systematic error of a model or an estimator, as in the bias--variance
decomposition (\cref{sec:biasvariance}). (2) The assumptions that a learning algorithm makes in order
to generalise, its inductive bias (\cref{sec:biasvariance}), or the language in which its models are
written, its language bias (Unit~4). (3) A systematic disadvantage for a group of people
(\cref{sec:fairness}).
\item[Calibration] (1) Agreement between predicted probabilities and observed frequencies
(\cref{sec:calibration}). (2) A \emph{calibration set}: held-out data used to fit a calibrator or the
threshold of conformal prediction (\cref{sec:calibration,sec:other}).
\item[Confidence] (1) The probability that a classifier gives to its predicted class
(\cref{sec:cautious}). (2) Of an association rule: the fraction of the examples that satisfy its
antecedent and also satisfy its consequent (\cref{sec:unsupervised}). (3) The level of a confidence
interval (\cref{sec:stats}).
\item[Counterfactual] (1) In causal inference: what would have happened under a different action
(\cref{sec:causality}). (2) A counterfactual explanation: the smallest change of the inputs that
changes the prediction (\cref{sec:examplebased}). (3) Counterfactual evaluation: measuring what a
deployed model changes, with a hold-out kept away from it (\cref{lst:holdout}) or with logged
decisions (\cref{lst:ips}).
\item[Coverage] (1) Of a rule: the fraction of the examples it applies to (\cref{sec:intmodels}).
(2) Of prediction sets: the fraction that contain the true label (\cref{sec:other}). (3) Of a cautious
classifier: the fraction of the examples on which it predicts (\cref{sec:cautious}). (4) Of a
recommender: the fraction of the catalogue it recommends (\cref{sec:receval}).
\item[Diversity] (1) Of an ensemble: how differently its members err (\cref{sec:crowds}). (2) Of a
list of recommendations: how different its items are (\cref{sec:receval}).
\item[Fidelity] (1) Of a surrogate model or an explanation: how well it reproduces the predictions
of the model it explains (\cref{sec:agnostic,sec:drawbacks}). (2) In multi-fidelity optimisation:
how closely a cheaper, approximate evaluation of a configuration follows the full one (Unit~4).
\item[Fitness] (1) In genetic programming: the score that decides which programs survive
(\cref{sec:gp}). (2) In process mining: how much of the observed behaviour a process model can replay
(\cref{sec:procmining}).
\item[Hold-out] (1) A test set kept apart from training (\cref{sec:protocols}). (2) A counterfactual
hold-out: cases or users kept away from a deployed model, so that its effect can be measured
(\cref{lst:holdout}).
\item[Lift] (1) In gains and lift charts: how many more positives a model finds than random
selection (\cref{sec:roc}). (2) Of an association rule: its confidence divided by the frequency of its
consequent (\cref{sec:unsupervised}). (3) In A/B testing: the relative improvement of a variant over
the control (\cref{sec:abtesting}).
\item[Metalearning] (1) In combinations of models: learning how to fuse the predictions of the base
models, as in stacking (\cref{sec:combining}). (2) Learning from previous experiments which algorithm or
configuration to use for a new problem (\cref{sec:metalearning}).
\item[Novelty] (1) Of a recommendation: how little known the recommended item is
(\cref{sec:receval}). (2) The novelty effect in experiments: users react to a change because it is new
(\cref{sec:abtesting}). (3) Novelty detection: recognising inputs of classes not seen in training
(\cref{sec:other}).
\item[Precision] (1) The fraction of positive predictions that are right (\cref{sec:crisp}). (2) The
number of bits used to store numbers, which quantisation reduces (\cref{sec:deployoptions}).
\item[Randomisation] (1) In ensembles: random choices while building the members, to make them
diverse (\cref{sec:ensembles}). (2) In experiments: random assignment of users to variants, which
makes the comparison causal (\cref{sec:abtesting}).
\item[Regret] (1) In metalearning: the performance lost by choosing an algorithm other than the best
one (\cref{sec:metalearning}). (2) In bandits: the reward lost over time by not always playing the
best arm (\cref{sec:bandits}).
\item[Support] (1) Of an association rule: the fraction of the transactions that contain all its
items (\cref{sec:unsupervised}). (2) Of a distribution: the region where it has data; importance
weighting fails where the training data have no support (\cref{sec:adapt}).
\item[Threshold] (1) The decision threshold on a score (\cref{sec:roc,sec:costs}). (2) The level at
which a monitor raises an alert (\cref{sec:detectshift}).
\item[Validation] (1) A validation set, used to select models and hyperparameters
(\cref{sec:protocols}). (2) Data validation: tests that the data meet their contract
(\cref{sec:infrastructure}).
\end{description}
